\subsection{Implementation and Experimental Setup}\label{sec:setup}
This chapter delineates the technical specifications and configurations utilized throughout the experimental phase of this research.
The following sections provide a comprehensive description of the hardware environment and the software architecture to facilitate the exact reproduction of the reported results. 
Thereafter, the implementation of InterCont will be described, while explaining the self-programmed loop.py.
The implementation of the benchmark methods will be described thereafter.


The experiments were conducted on a workstation PC equipped with an Intel Core i9-12900K processor operating at 3.2 GHz and an NVIDIA GeForce RTX 3090 GPU. 
The system features 64 GB of physical RAM and 24 GB of video memory (VRAM), providing sufficient computational capacity. 
The software environment was managed under the Windows 11 Enterprise operating system, specifically version 25H2. 


For the implementation of the algorithms and the execution of the experimental pipeline, the PyCharm 2025.3.3 IDE was employed. 
Although more recent versions of the Python interpreter exist, Python 3.12 \citep{pythonpsf} was selected as the foundational language to ensure full compatibility with the required library ecosystem, as several critical packages had not yet been updated for newer releases. 
To maintain strict control over software dependencies and ensure environment stability, a virtual environment was utilized. 
All specific library versions and dependencies are documented within a provided requirements.txt file.


\subsubsection*{InterCont}
The authors of InterCont provided a comprehensive reproduction package alongside their publication, which enabled a seamless replication of the original results. 
However, the provided codebase did not include implementations for k-NN. 
Consequently, k-NN and MCD were independently implemented for this study. 


The reproduction package is structured into four distinct Python modules: main.py, train.py, data\_loader.py, and helper\_functions.py. 
The main.py script serves as the primary orchestrator for the experimental pipeline, managing configuration parameters, initiating data loading procedures, and executing the training sequence. 
The adjustable parameters include the dataset selection, an option for an accelerated version of the algorithm, and the batch size. 
The default batch size is set to 3000, which, given the dimensions of the datasets utilized in this research, effectively results in a full-batch gradient descent.
This parameter was adjusted to 129, to ensure full batch learning. 


The dataset utilized for the experiments can be specified via command-line arguments.
However, since the provided reproduction package includes the entire ODDS repository and the default value selected the Wine dataset, this parameter remained unchanged. 
Furthermore, the configuration includes an option to enable an accelerated version of the algorithm. 
This feature was not activated, as the primary dataset is of a manageable size that does not necessitate the optimizations offered by the fast implementation. 
To ensure computational efficiency, the authors included a logic sequence at the conclusion of the execution script to prioritize the use of the GPU. 


At the initialization of this module, a data pipeline is established through a specialized data wrapper, which can be found in the data\_loader.py. 
For this implementation the function def build\_train\_test\_generic\_matfile is used.
This component determines the structural dimensions of the dataset and defines the  for individual row retrieval.


This function is essential for the implementation of the benchmark methods. 
To facilitate the subsequent evaluation of calculated anomaly scores, the wrapper preserves the original index identifiers. 
This architectural choice ensures that the machine learning algorithm can accurately map each anomaly score back to its corresponding observation in the source data.


To introduce stochasticity into the shuffle process, the authors utilized the PyTorch \citep{paszke2019pytorch} function torch.randperm().
This function divides the normal data into halves and shuffles the compositions and grading those halves into testing and training dataset.
The testing dataset always includes all outliers.  
To ensure the stability and reliability of the results, a shuffling procedure is integrated into the pipeline. 
The torch.randperm() function relies on a pseudo-random number generator.
Since no specific generator argument was specified, the default engine was employed. 


As the reproduction package did not include a native looping mechanism for these iterations, a custom loop.py was implemented to manage the process. 
Finally, performance metrics were derived in accordance with the logic defined within the helper\_functions module.
Consistent with the original study, the experiments were executed over 500 random splits. 


While the original codebase was designed to accommodate a vast array of datasets, including specialized logic within the data loader for .mat files, many components proved redundant for the specific dataset analyzed in this research. 
This code was retained to preserve the integrity of the reproduction package, as redundant logic does not influence the experimental outcomes. 
Notably, several hyperparameters within the original implementation were optimized for larger, higher-dimensional datasets. 
With exception of the batch size, none of the hyperparameters were changed.
Furthermore, while the original authors restricted their evaluation to F1-Score and ROC AUC metrics, this study extends the analysis to include AUPRC. 
Minor adjustments were made to the source code to facilitate this additional metric.


\subsubsection*{k-NN}
The k-NN algorithm was implemented in knn\_benchmark.py by using the PyOD library \citep{zhao2019pyod}. 
To facilitate the experimental procedure and derive relevant performance metrics, components from the Scikit-learn library \citep{pedregosa2011scikit} were also employed. 
Especially the F1-Score provided by Scikit-learn was used to estimate this metric. 
Nevertheless, by using the known outlier score, the use of both methods yield to the same result.


A five-fold cross-validation procedure was implemented in knn\_cv.py.
This was facilitated by the KFold function from the Scikit-learn library. 
$K = 5$ is the optimal value for $k$, which corresponds with the choice in the paper of the authors. 
Given the limited number of observations in the dataset, a five-fold approach was selected to maintain a sufficient volume of data in both the training and testing subsets for each iteration. 
The shuffle parameter $random\_state$ was enabled to ensure that the observations were randomized prior to partitioning and was set to the default 42. 


To ensure a rigorous and consistent evaluation, the same framework utilizing the torch.randperm () function, was integrated into the $k$-NN implementation. 
In this configuration, the $k$-NN algorithm establishes a representation based on the normal observations within the training set. 
Subsequently, the method calculates the distances between each test observation and its $k$ nearest neighbors from the training data. 
The anomaly score for each test point is derived from the maximum distance among these $k$ nearest neighbors. 
The dataset includes 119 normal datapoints, which can not be separated equally into training and testing.
To align the model with the specific characteristics of the test dataset, the number of outliers are hardcoded into the script.


The $k$-NN implementation utilizes the default configuration for distance metrics, specifically selecting the largest distance to the $k$-nearest neighbors. 
This calculation employs the Minkowski distance with the parameter $p=2$, which is mathematically equivalent to the Euclidean distance.
The RobustScaler from Scikit-learn was implemented.

\subsubsection*{FAST-MCD}
The FAST-MCD algorithm is implemented in mcd\_benchmark.py utilizing the PyOD library. 
As an unsupervised method, FAST-MCD does not support a traditional separation into training and testing phases. 
If a split-sample approach were attempted, the algorithm would collapse.
Specifically, if $n=h$, the resulting  $\mu_MCD$ and $\sum_MCD$ matrix would be equivalent to the mean and covariance matrix of the training dataset \citep[p. 217]{rousseeuw1999fast}. 
To ensure experimental comparability, the FAST-MCD algorithm is applied directly to the test dataset. 
The training dataset is generated using the same shuffle function employed by the other evaluated methods. 
Observations that exhibit the 10 highest Mahalanobis distances relative to the estimated mean and covariance matrix are classified as outliers, adhering to a top\_k=10 criterion. 


To ensure parity with the implementation of InterCont, this cycle was repeated 500 times. 
The final performance metrics represent the mean values across these iterations, supplemented by the calculation of the standard deviation. 
The random\_state parameter was fixed at 42 to guarantee reproducibility, because FAST-MCD relies on selecting random initial subsets before performing subsequent C-steps.
















